\chapter{The laws}
\label{chap:laws}

The essay states three laws of computation, by analogy with thermodynamics. The zeroth and the
first are theorems; the second is a named hypothesis in two strengths.

\section{The zeroth law}

\begin{theorem}[Zeroth law]
  \label{thm:lawZero}
  \lean{Sagitta.Arrow.Model.lawZero}
  \leanok
  \uses{def:model}
  In every model, reductions are transitive: if $A \red B$ and $B \red C$ then $A \red C$.
  \own{} (the name and the reading as a law; transitivity itself is standard,
  \cite{AroraBarak}, Chapter~2).
\end{theorem}

\begin{proof}
  \leanok
  It is the closure property ``reductions compose'' of the model.
\end{proof}

\begin{corollary}[Zeroth law for Turing machines]
  \label{cor:lawZero-tm}
  \lean{Sagitta.Arrow.TM.lawZero, Cslib.Complexity.PolyTimeReducible.trans}
  \leanok
  \uses{thm:lawZero, def:tmModel, def:reduction}
  Polynomial-time many-one reductions between encoded problems compose.
\end{corollary}

\begin{proof}
  \leanok
  Theorem~\ref{thm:lawZero} applied to the Turing-machine model; the field is cslib-frontier's
  \texttt{PolyTimeReducible.trans}.
\end{proof}

\begin{lemma}[Completeness travels along reductions]
  \label{lem:npcomplete-reduces}
  \lean{Sagitta.Arrow.Model.NPComplete.of_reduces, Sagitta.Arrow.Model.NPComplete.InP_iff}
  \leanok
  \uses{thm:lawZero, def:model-classes}
  If $B$ is $\cNP$-complete, $C \in \cNP$ and $B \red C$, then $C$ is $\cNP$-complete. An
  $\cNP$-complete $B$ is in $\cP$ if and only if $\cP = \cNP$.
\end{lemma}

\begin{proof}
  \leanok
  Compose the reductions; for the second part use the closure of $\cP$ under reductions.
\end{proof}

\section{The first law}

A \emph{query tree} over a finite set of options $\alpha$ asks, at each node, whether an option is
good, and answers yes or no at the leaves; its cost on an oracle $f : \alpha \to \{0,1\}$ is the
number of questions on the path that $f$ selects.

\begin{definition}[Query trees that decide existence]
  \label{def:querytree}
  \lean{Sagitta.Arrow.QueryTree, Sagitta.Arrow.QueryTree.cost, Sagitta.Arrow.QueryTree.DecidesExists}
  \leanok
  A query tree $t$ \emph{decides existence} when, for every oracle $f$, it answers yes exactly when
  $f(i) = 1$ for some option $i$.
\end{definition}

\begin{lemma}[The adversary]
  \label{lem:adversary}
  \lean{Sagitta.Arrow.QueryTree.adversary, Sagitta.Arrow.QueryTree.cost_adversary}
  \leanok
  \uses{def:querytree}
  A query tree that decides existence asks, on the all-zero oracle, about every option; so its cost
  there is at least $|\alpha|$.
  \source{the adversary argument for unstructured search; see \cite{AroraBarak}, Chapter~12
  (decision trees).}
\end{lemma}

\begin{proof}
  \leanok
  If an option $i$ is never asked on the all-zero oracle, the oracle that is $1$ only at $i$ follows
  the same path and gets the same answer, which is wrong for one of the two.
\end{proof}

\begin{theorem}[First law: work is conserved in a black box]
  \label{thm:conservation}
  \lean{Sagitta.Arrow.DependsOnlyOnAsked, Sagitta.Arrow.conservation}
  \leanok
  \uses{lem:adversary}
  Let a filling phase build a table $\mathrm{fill}(f)$ asking the options $\mathrm{asked}(f)$, and let
  the table depend only on the answers to those options (the system is closed). If a query phase
  $Q(\mathrm{fill}(f))$, reading the table for free, decides whether some option is good for every
  oracle $f$, then on the all-zero oracle
  \[
    |\alpha| \le |\mathrm{asked}(0)| + \mathrm{cost}(Q(\mathrm{fill}(0)), 0).
  \]
  What is hard is neither created nor destroyed: it is only moved between the phases.
  \own{} (the statement as a conservation law; the proof is the adversary argument).
\end{theorem}

\begin{proof}
  \leanok
  \uses{lem:adversary}
  Every option is asked by one of the two phases on the all-zero oracle; otherwise flip it and both
  phases behave the same.
\end{proof}

\begin{proposition}[Control: information from outside breaks the first law]
  \label{prop:free-table}
  \lean{Sagitta.Arrow.free_table, Sagitta.Arrow.free_table_not_closed, Sagitta.Arrow.free_table_breaks}
  \leanok
  \uses{thm:conservation}
  If the table is the oracle itself, the query phase decides at cost $0$ and the filling phase asks
  nothing; the table then does not depend only on what was asked, and the inequality of
  Theorem~\ref{thm:conservation} fails.
  \own{}
\end{proposition}

\begin{proof}
  \leanok
  Direct computation.
\end{proof}

\section{The second law}

\begin{definition}[The arrow, law 2.1]
  \label{def:arrow}
  \lean{Sagitta.Arrow.Model.Arrow}
  \leanok
  \uses{def:model-classes}
  The \emph{arrow} of a model is $\neg(\cP = \cNP)$, the worst-case version. It is a
  \textbf{hypothesis}: a definition of type \texttt{Prop}, taken as an explicit argument by every
  theorem that uses it, never an axiom.
  \own{} (the reading as a second law; the statement is $\cP \neq \cNP$).
\end{definition}

\begin{theorem}[Kelvin's form]
  \label{thm:kelvin}
  \lean{Sagitta.Arrow.Model.kelvin}
  \leanok
  \uses{def:arrow, thm:searchEasy-iff}
  The arrow holds if and only if search is not easy: there is no efficient process whose only
  effect is to convert every ``yes'' into its witness.
  \own{}
\end{theorem}

\begin{proof}
  \leanok
  \uses{thm:searchEasy-iff}
  The localization, part (3): search is easy exactly when $\cP = \cNP$.
\end{proof}

\begin{definition}[Two models]
  \label{def:two-models}
  \lean{Sagitta.Arrow.collapsedModel, Sagitta.Arrow.separatedModel}
  \leanok
  \uses{def:model}
  The \emph{collapsed} model has one problem, in $\cP$, and every search problem easy; the
  \emph{separated} model satisfies the closure properties with $\cP \neq \cNP$.
  \own{}
\end{definition}

\begin{theorem}[The arrow is independent of the closure properties]
  \label{thm:arrow-independent}
  \lean{Sagitta.Arrow.arrow_independent}
  \leanok
  \uses{def:arrow, def:two-models}
  The arrow fails in the collapsed model and holds in the separated model.
  \own{}
\end{theorem}

\begin{proof}
  \leanok
  Compute in both models.
\end{proof}

\begin{definition}[The arrow for Turing machines]
  \label{def:arrow-tm}
  \lean{Sagitta.Arrow.TM.Arrow}
  \leanok
  \uses{def:arrow, def:tmModel}
  The arrow of the Turing-machine model: for some encoded type, $\cNP \not\subseteq \cP$. A
  hypothesis, never an axiom.
\end{definition}

\begin{theorem}[The arrow, unfolded]
  \label{thm:arrow-iff-tm}
  \lean{Sagitta.Arrow.TM.arrow_iff, Sagitta.Arrow.TM.arrow_iff_exists, Sagitta.Arrow.TM.arrow_iff_SAT, Sagitta.Arrow.TM.kelvin}
  \leanok
  \uses{def:arrow-tm, thm:kelvin, thm:cook-levin, thm:oneBit}
  For Turing machines the arrow is equivalent to each of: not every $\cNP$ class is contained in
  $\cP$; $\cP \neq \cNP$ for some encoding; $\SAT \notin \cP$; and (Kelvin's form) not every $\cFNP$
  search problem is in $\cFP$.
\end{theorem}

\begin{proof}
  \leanok
  \uses{thm:cook-levin, thm:oneBit}
  The third equivalence is the Cook--Levin theorem; the fourth is cslib-frontier's
  search-to-decision theorem.
\end{proof}

\begin{definition}[Law 2.2]
  \label{def:law22}
  \lean{Sagitta.Arrow.TM.InTIME, Sagitta.Arrow.TM.subexpBound, Sagitta.Arrow.TM.Law22}
  \leanok
  \uses{def:tmModel, def:threesat}
  \emph{Law 2.2}: for every $\varepsilon > 0$, $\TSAT \notin \TIME(2^{n^{1-\varepsilon}})$, where
  $\TIME(T)$ allows $C \cdot T(n) + C$ steps for a constant $C$ and $n$ is the length of the encoded
  formula. A hypothesis, never an axiom.
  \source{the weak exponential-time hypothesis, shown equivalent to Stearns and Hunt's satisfiability
  hypothesis SH in \cite{ETHdead} (as cited in the Lean source, Theorem~2.4).}
  \fidelity{$n$ is measured in bits of the CNF encoding (the size form of SH), not in variables. It is
  equivalent to the weak ETH ($n$ variables) only by Theorem~2.4 of \cite{ETHdead}, which is not
  formalized here.}
\end{definition}

\begin{theorem}[Law 2.2 implies law 2.1]
  \label{thm:law22-arrow}
  \lean{Sagitta.Arrow.TM.Law22.arrow, Sagitta.Arrow.TM.poly_le_subexp}
  \leanok
  \uses{def:law22, def:arrow-tm, def:threesat}
  Law 2.2 implies the arrow.
\end{theorem}

\begin{proof}
  \leanok
  If $\cP = \cNP$ then $\TSAT$ is decided in polynomial time, hence in time
  $C \cdot 2^{n^{1/2}} + C$, which law 2.2 forbids with $\varepsilon = 1/2$.
\end{proof}
